% Auto-generated from meeting_2568-11-15_interview-committee.md (กบ.วช. format v2)
\documentclass{meetingminutes}
\meetingcommittee{คณะกรรมการบริหารหลักสูตร วิศวกรรมคอมพิวเตอร์และปัญญาประดิษฐ์ (บัณฑิตศึกษา)}
\meetingnumber{๑}{๒๕๖๘}
\meetingdate{วันที่ ๑๕ พฤศจิกายน พ.ศ. ๒๕๖๘}
\meetingplace{ห้องประชุมคณะเทคโนโลยีอุตสาหกรรม}
\meetingstart{๐๘.๓๐ น.}
\meetingend{๑๕.๒๐ น.}

\begin{document}
\maketitle

\psection{ผู้มาประชุม}
\begin{participants}
\end{participants}

\psection{ผู้ไม่มาประชุม}
\noindent ไม่มี\par

\startmeeting
\openingchair{ผู้ช่วยศาสตราจารย์ ดร.พิศิษฐ์ นาคใจ}{กล่าวเปิดการประชุมและดำเนินการประชุมตามระเบียบวาระ ดังนี้}

\agenda{๑}{สรุปจำนวนผู้สมัคร}
ฝ่ายเลขานุการรายงานว่ามีผู้สมัครเข้าศึกษาในหลักสูตร วศ.ม. วิศวกรรมคอมพิวเตอร์และปัญญาประดิษฐ์ ปีการศึกษา 2568 (ภาคเรียนที่ 2) จำนวนทั้งสิ้น 6 คน ซึ่งผ่านการตรวจสอบคุณสมบัติเบื้องต้นครบทุกคน และได้รับการนัดหมายสอบสัมภาษณ์วันนี้


\agenda{๒}{ดำเนินการสอบสัมภาษณ์}
คณะกรรมการดำเนินการสอบสัมภาษณ์ผู้สมัครทั้ง 6 คน ในช่วงเวลา 09.00 – 14.00 น. โดยประเมินตามเกณฑ์ที่กำหนด: 1. ความรู้พื้นฐานด้านวิศวกรรมคอมพิวเตอร์/วิทยาการคอมพิวเตอร์/AI 2. แรงจูงใจและเป้าหมายในการศึกษาต่อระดับมหาบัณฑิต 3. แผนการวิจัยหรือสารนิพนธ์เบื้องต้น 4. ความสามารถในการสื่อสารและการแสดงความคิดเห็น


\agenda{๓}{สรุปผลการสัมภาษณ์และมติรับเข้าศึกษา}
คณะกรรมการร่วมพิจารณาผลการสัมภาษณ์ผู้สมัครทั้ง 6 คน

\begin{longtable}{|p{2.5cm}|p{4cm}|p{6cm}|}
  \hline
  \textbf{ลำดับ} & \textbf{สถานะ} & \textbf{หมายเหตุ} \\
  \hline
  1–6 & สอบสัมภาษณ์ครบ & ผู้สมัครทั้ง 6 คนผ่านการสัมภาษณ์ \\
  \hline
\end{longtable}

\resolution{รับผู้สมัครทั้ง 6 คนเข้าศึกษา และดำเนินการแจ้งผลและนัดรายงานตัวตามกำหนด}

\agenda{๔}{กำหนดการรายงานตัวและขึ้นทะเบียนนักศึกษา}
ที่ประชุมกำหนดการรายงานตัว: - กำหนดรายงานตัว: ต้นธันวาคม 2568 - นักศึกษาที่รายงานตัวจริง: 4 คน (2 คน ไม่มารายงานตัวตามกำหนด)

นักศึกษาที่รายงานตัว (รุ่นที่ 1): 1. นายเกริกเกียรติ บุญไทย (รหัส 68556680101) — หัวข้อวิจัย: Deep Learning ตรวจจับผู้ไม่สวมหมวกนิรภัย 2. นางสาวพอใจ สุขหา (รหัส 6888668106) — หัวข้อวิจัย: IoT ตู้อบแห้งข้าวแต๋นอัจฉริยะ 3. นายจาตุรงค์ ทรงภักดี — หัวข้อวิจัย: AI จัดการสถานีชาร์จ EV 4. นายจิรกิตติ์ วราภรณ์ (รหัส 68556680104) — หัวข้อวิจัย: Remote Sensing ระบุพื้นที่ปลูกอ้อย

\resolution{รับทราบจำนวนนักศึกษาที่รายงานตัวจริง 4 คน และดำเนินการขึ้นทะเบียนต่อไป}

\adjournmeeting

\begin{sigblock}
\typist{ผู้ช่วยศาสตราจารย์ ดร.กาญจนา ดาวเด่น}
\reviewer{ผู้ช่วยศาสตราจารย์ ดร.พิศิษฐ์ นาคใจ}{กรรมการ}
\chairsig{ผู้ช่วยศาสตราจารย์ ดร.พิศิษฐ์ นาคใจ}{ประธานการประชุม}
\end{sigblock}

\end{document}
